% \setcounter{chapter}{\Session-2}
\renewcommand{\chaptername}{سوال}
\chapter{حل تمرین دوم: مدل‌های مولد و کاربردها}

در این بخش پاسخ‌های مفصل به سؤال‌های تمرین دوم ارائه می‌شود. ساختار هر پاسخ شامل «صورت سؤال»، «رویکرد پیشنهادی»، «جزئیات پیاده‌سازی»، «نتایج مورد انتظار» و «تحلیل و جمع‌بندی» می‌باشد.

\section{سؤال ۱ — توضیح نظریهٔ جریان‌های نرمال‌سازی‌شده (MAF)}
\subsection{صورت سؤال}
توصیف کنید جریان‌های نرمال‌سازی‌شده چه هستند، اصول ریاضی پشت آن‌ها را شرح دهید و نشان دهید چگونه MAF از فاکتوراسیون اتورگرسیو برای مدل‌سازی چگالی استفاده می‌کند.

\subsection{پاسخ}
جریان‌های نرمال‌سازی‌شده مجموعه‌ای از نگاشت‌های بازگشت‌پذیر $f_{\theta}$ هستند که یک توزیع سادهٔ پایه مانند گاوسی را به توزیع دادهٔ هدف نگاشت می‌کنند؛ با استفاده از تغییر متغیرها می‌توان لاگ‌احتمال را به‌صورت دقیق محاسبه کرد:
\begin{equation}
\log p_{\theta}(x) = \log p_Z(f_{\theta}(x)) + \log\left|\det \frac{\partial f_{\theta}(x)}{\partial x}\right|.
\end{equation}
در MAF از فاکتورسازی اتورگرسیو استفاده می‌شود تا چگالی مشترک به حاصل‌ضرب شرطی‌ها شکسته شود:
\begin{equation}
p(x)=\prod_{d=1}^D p(x_d\mid x_{<d}).
\end{equation}
هر شرطی با یک شبکهٔ عصبی پارامتریزه می‌شود و با ماسک‌گذاری (MADE) وابستگی‌های مجاز تضمین می‌گردد. در نتیجه می‌توان پارامترهای مقیاس و میانگین را به‌صورت شرطی محاسبه کرد و نگاشت‌های ساده‌ای با ژاکوبی مثلثی ساخت که دترمینان آن‌ها حاصل‌ضرب مؤلفه‌های قطری است؛ این امر محاسبهٔ ژاکوبی را با پیچیدگی خطی ممکن می‌سازد.

\section{سؤال ۲ — طراحی و آموزش CycleGAN}
\subsection{صورت سؤال}
معماری CycleGAN را تشریح کنید و توضیح دهید چگونه ضررهای چرخه‌ای و هویتی به حفظ ساختار تصاویر کمک می‌کنند. همچنین مراحل آموزش و نکات پایدارسازی را فهرست کنید.

\subsection{پاسخ}
معماری شامل دو ژنراتور ($G: X\to Y$، $F: Y\to X$) و دو افتراق‌دهنده ($D_X, D_Y$) است. ژنراتورها معمولاً مبتنی بر بلوک‌های ResNet و افتراق‌دهنده‌ها از ساختار PatchGAN استفاده می‌کنند تا تمرکز بر بافت‌های محلی حفظ شود. ضرر چرخه‌ای
\(\mathcal{L}_{cyc}=\mathbb{E}_{x}[\|F(G(x))-x\|_1]+\mathbb{E}_{y}[\|G(F(y))-y\|_1]\)
باعث می‌شود نگاشت‌ها تقریباً معکوس یکدیگر باشند و از جابجایی شدید ساختار جلوگیری گردد. ضرر هویتی نیز برای حفظ رنگ و روشنایی بهره‌برده می‌شود:
\(\mathcal{L}_{id}=\mathbb{E}_{y}[\|G(y)-y\|_1]+\mathbb{E}_{x}[\|F(x)-x\|_1]\).

نکات پایدارسازی: استفاده از LSGAN یا نرم‌سازی برچسب‌ها، بافر بازپخش نمونه‌های جعلی، زمان‌بندی نرخ یادگیری با کاهش خطی بعد از مجموعه‌ای از اپک‌ها، و به‌کارگیری نرمال‌سازی نمونه (Instance Norm) در ژنراتور.

\section{سؤال ۳ — تشخیص ناهنجاری با جریان‌ها}
\subsection{صورت سؤال}
شرح دهید چگونه می‌توان از مدل‌های مبتنی بر لاگ‌احتمال (مثل MAF) برای تشخیص ناهنجاری استفاده کرد و محدودیت‌های عملی این رویکرد را بیان کنید.

\subsection{پاسخ}
برای تشخیص ناهنجاری از نمرهٔ منفی لاگ‌احتمال (NLL) به‌عنوان معیار استفاده می‌شود؛ نمونه‌هایی که NLL زیادی دارند به‌عنوان ناهنجاری شناخته می‌شوند. محدودیت‌ها شامل: کالیبراسیون نامناسب احتمال برای تصاویر پیچیده، حساسیت به تغییرات سطحی مانند روشنایی، و نیاز به دادهٔ آموزش تقریباً کامل از وضعیت «نرمال». برای بهبود می‌توان از ویژگی‌های استخراج‌شدهٔ CNN و یا ترکیب با روش‌های مبتنی بر بازسازی استفاده کرد.

\section{سؤال ۴ — جزئیات پیاده‌سازی و تنظیمات تجربی}
\subsection{صورت سؤال}
جدول تنظیمات هایپرپارامترها، پیش‌پردازش داده‌ها و روش‌های ارزیابی را شرح دهید.

\subsection{پاسخ}
خلاصه‌ای از تنظیمات (نمونه):
\begin{itemize}
	\item MAF: بلوک‌ها = 7، عرض پنهان = 512، lr = $1\times10^{-4}$، batch = 3، اپک = 100.
	\item CycleGAN: بلوک‌های ResNet = 9، lr = $2\times10^{-4}$، batch = 1، epochs = 200 (خطی کاهش lr از اپک 100).
	\item پیش‌پردازش: تغییر اندازه، نرمال‌سازی کانال‌ها، برای جریان‌ها فلت‌کردن تصویر پس از استانداردسازی.
	\item ارزیابی: AUROC برای تشخیص ناهنجاری، ارزیابی‌های کیفی و در صورت نیاز FID برای ترجمهٔ تصویر.
\end{itemize}

\section{سؤال ۵ — تحلیل نتایج و پیشنهادات برای بهبود}
\subsection{صورت سؤال}
نتایج تجربی را تحلیل کنید و پیشنهاداتی برای بهبود مدل‌ها ارائه دهید.

\subsection{پاسخ}
تحلیل کلی: MAF برای تشخیص ناهنجاری مزیت محاسبهٔ صریح احتمال را دارد اما در تصاویر با ساختار مکانی قوی ضعف دارد؛ CycleGAN ترجمه‌های واقع‌گرایانه تولید می‌کند اما در تغییرات هندسی پیچیده مشکل دارد. پیشنهادها:
\begin{itemize}
	\item استفاده از جریان‌های کانولوشنی (Glow، Residual Flows) برای بهره‌گیری از ساختار مکانی.
	\item تلفیق لاگ‌احتمال با ویژگی‌های بلندمدت استخراج‌شده توسط شبکه‌های پیش‌آموزش‌دیده.
	\item اضافه کردن معیارهای ساختاری (مثل SSIM) به تابع هزینهٔ CycleGAN.
\end{itemize}

\section{جمع‌بندی}
پاسخ‌های بالا به‌صورت کامل و رسمی ارائه شدند و می‌توانند مستقیماً در گزارش اصلی مرجع یا پیوست قرار گیرند. در صورت تمایل من می‌توانم همین ساختار را برای `session1.tex` نیز اعمال یا هر پاسخ را بر اساس صورت سؤال‌های دقیق شما بازنویسی کنم.
